%---------------------------------------------------------------------------------------

\chapter[Considerações Finais]{Considerações Finais}\label{ch:consideracoes-finais}

%---------------------------------------------------------------------------------------

% Lista de abreviaturas e de siglas

%---------------------------------------------------------------------------------------

% Lista de simbolos

%---------------------------------------------------------------------------------------

Este trabalho apresentou um miniaplicativo que implementa o método de Lattice–Boltzmann para escoamentos compressíveis, validado no problema do tubo de choque de Sod e avaliado quanto à escalabilidade em CPUs multinúcleo.
Os resultados demonstram (i) consistência numérica nos regimes com descontinuidades, com taxa efetiva próxima da primeira ordem em normas globais, e (ii) bom aproveitamento de paralelismo de memória compartilhada para tamanhos de problema moderados a grandes, com eficiência elevada até o limite imposto por largura de banda e sincronizações.

Do ponto de vista científico, a acurácia observada na razão de pressões pós-interação e a tendência de redução do EQM reforçam a viabilidade do LBM compressível como alternativa competitiva para problemas hiperbólicos com ondas de choque, preservando simplicidade algorítmica e regularidade de acesso à memória.
Esse equilíbrio entre fidelidade física e eficiência computacional é particularmente atraente para estudos paramétricos e acoplamentos multiescala, onde o custo por simulação é crítico.

Em termos de impacto acadêmico, o estudo fornece uma linha de base reprodutível para comparações entre esquemas de relaxação, discretizações em velocidade e estratégias de filtragem ou dissipação controlada.
Além disso, apresenta um perfil de desempenho que explicita gargalos práticos, tais como largura de banda, afinidade e sincronizações, orientando pesquisas em otimização de \textit{kernels} LBM no contexto de arquiteturas heterogêneas.
Destaca-se também o desenvolvimento de um artefato experimental útil para ensino e extensão, permitindo a estudantes e pesquisadores a reprodução de curvas de convergência e escalabilidade com esforço modesto.

Como trabalhos futuros, propõe-se um roteiro de evolução tecnológica e experimental estruturado em múltiplas frentes.
Primeiramente, no âmbito de portabilidade de desempenho (\textit{performance portability}), sugere-se a reestruturação dos \textit{kernels} de colisão e \textit{streaming} utilizando a biblioteca Kokkos.
Essa abordagem permitirá explorar espaços de execução em OpenMP, CUDA ou HIP, além de \textit{layouts} de dados compatíveis com coalescência e vetorização.
Em paralelo, a criação de uma variante com RAJA permitirá a avaliação de abstrações alternativas de políticas de execução e particionamento (\textit{tiling}), comparando-se sobrecusto computacional e expressividade algorítmica.

Na frente de paralelismo distribuído, propõe-se a implementação de decomposição de domínio multidimensional com MPI, utilizando camadas fantasmas (\textit{ghost layers}) e comunicação não bloqueante, aliadas à agregação de mensagens para mitigação de latência.
Tal arquitetura permitirá a sobreposição de comunicação e computação através de processos assíncronos e estratégias de balanceamento de carga para partições desiguais em sistemas heterogêneos.

Para uma avaliação rigorosa, faz-se necessária uma campanha extensiva de \textit{benchmarks} em ambiente de computação de alta performance.
Isto inclui ensaios de escalabilidade forte e fraca em \textit{clusters} equipados com aceleradores de GPU e nós robustos de CPU, variando o refinamento espacial e o número de passos de tempo.
Nestas campanhas, devem ser mensurados o \textit{throughput} (em MLUPS), a eficiência paralela, a correlação através de modelos \textit{roofline} e o custo energético.
Complementarmente, estudos de sensibilidade a parâmetros físicos e numéricos, como filtros estabilizadores, permitirão o mapeamento do nível da estabilidade versus o custo computacional incorrido.

Adicionalmente, a engenharia de desempenho do código pode ser aprimorada mediante a adoção de \textit{layout} SoA com blocagem espacial, busca antecipada (\textit{prefetching}) e alinhamento explícito, bem como fusão de \textit{kernels} para diminuição da pressão sobre a memória.
A realização de validações físicas com problemas com descontinuidades mais severas, incluindo choques de Lax e de Shu-Osher, confrontando-se os resultados numéricos com soluções exatas ou de referência, confirmará a resiliência numérica do método LBM compressível nas condições operacionais complexas.

Concluímos que o miniaplicativo alcança um ponto de maturidade que justifica a migração para um arcabouço de portabilidade e sua avaliação em infraestruturas de HPC. A expectativa é sustentar eficiência próxima ao teto de largura de banda em CPUs modernas e explorar aceleração substancial em GPUs, preservando a robustez numérica nas regimes com descontinuidades.
Essa trajetória posiciona o LBM compressível como ferramenta de vanguarda para investigações acadêmicas e aplicações industriais que demandam alto rendimento computacional.
